\documentclass[10pt,a4paper]{amsart}
\usepackage[margin=19mm]{geometry}
\usepackage{amsmath,amssymb,mathtools}
\usepackage[scr=rsfs,cal=euler]{mathalfa}
\usepackage{xfrac}
\usepackage[unicode,hidelinks,bookmarks=false]{hyperref}
\newcommand{\ie}{i.e.\ }
\newcommand{\cf}{cf.\ }
\newcommand{\resp}{respectively\ }
\newcommand{\Subsection}{\subsection}
\providecommand{\nobreakdash}{\nobreak\hbox{-}}
\setlength{\emergencystretch}{2em}
\multlinegap0pt
\usepackage{xcolor}
\usepackage{algorithm}\usepackage{algpseudocode}
\newtheorem{lemma}{Lemma}
\title{Information-Geometric Optimization on Spheres}
\author{Vladimir Ja\' cimovi\'c}
\date{Source: arXiv:2606.07588v1 (CC BY 4.0)}
\begin{document}
\maketitle

\noindent\emph{Source and licence.} Information-Geometric Optimization on Spheres. Version arXiv:2606.07588v1 (CC BY 4.0), distributed by arXiv under the Creative Commons Attribution 4.0 International licence, http://creativecommons.org/licenses/by/4.0/, as recorded in the arXiv licence metadata for this exact version and shown on the abstract page https://arxiv.org/abs/2606.07588v1. Full licence notice: \texttt{licenses/arxiv-2606.07588-CC-BY-4.0.txt}.\par\smallskip

\noindent\textit{Editorial scope.} Counted unit: the article's lemma Moebius\_lemma (source0.tex lines 975--987). The article's complete original proof of it is delivered with it (source0.tex lines 989--1011, one \texttt{proof} environment of 23 lines). No mathematics is written, repaired or re-derived: the statement and the proof are the article's own lines, in the article's own notation, and no retained line is substituted or re-flowed.  No further source-local context is needed: every cross-reference of every retained line resolves inside the two delivered spans.  Nothing was omitted from any retained span, so the package carries no omission record.  No retained line contains a citation, so the wrapper carries no bibliography and no bibliography entry.  No retained line contains a figure environment, an image-inclusion command or drawing code, so no image is delivered and the package declares no asset dependency.  Editorial formatting only, in addition: an AMS \texttt{amsart} preamble that restates the article's own declarations and macros used by the retained lines, namely the environments \texttt{lemma} on separate counters, together with the standard packages those lines need.  The retained environments print in this document as Lemma 1 in order of appearance, whereas the article prints the same items with the numbering of its own full document.  The complete native source of the article as distributed in the arXiv e-print for this exact version is bundled unchanged as \texttt{sources/202265/source0.tex}.
\begin{lemma}
Suppose that vector $v \in {\mathbb C}^m$ satisfies
\begin{equation}
\label{Moebius_lemma}
v = R \frac{\beta - w}{1 - w^\dagger \beta} \quad \mbox{  with } \; R \in U(d).
\end{equation}
Then
\begin{equation}
\label{identity1}
\frac{1}{1 - w^\dagger \beta} = \frac{1 + w^\dagger R^\dagger v}{1 - \| w\|^2} \quad \mbox{and} \quad 
\beta = \frac{R^\dagger v + w}{1 + w^\dagger R^\dagger v}.
\end{equation}
\end{lemma}

\begin{proof}
Multiply \eqref{Moebius_lemma} by $R^\dagger$ to get
\begin{equation}
\label{expr_lemma}
R^\dagger v (1 - w^\dagger \beta) = \beta - w \implies \beta = w + R^\dagger v (1 - w^\dagger \beta).
\end{equation}
Multiply the last equality by $w^\dagger$ to get
$$
w^\dagger \beta = \| w \|^2 + w^\dagger R^\dagger v (1 - w^\dagger \beta).
$$
Substitute the last equality from one to get
$$
1 - w^\dagger \beta = 1 - \|w\|^2 - w^\dagger R^\dagger v (1 - w^\dagger \beta) \implies (1 - w^\dagger \beta) (1 + w^\dagger R^\dagger v) = 1 - \|w\|^2.
$$
Rearranging multipliers in the last equality yields the first (left-side) expression in \eqref{identity1}.

In order to derive the second equality in \eqref{identity1} substitute the first equality into \eqref{Moebius_lemma}. This yields:
$$
v = R (\beta - w) \frac{1 + w^\dagger R^\dagger v}{1 - \| w \|^2}.
$$
Now, it is easy to solve the last expression w. r. to $\beta$ to get the right-hand equality in \eqref{identity1}.

\end{proof}

\end{document}
